\chapter{SIGIL Pathways}
\label{chap:SIGIL_Pathways}

\section{Purpose}

SIGIL Pathways define how human-in-the-loop (HITL) operators use symbolic keys,
rituals, or ``sigils'' to gate access to critical CME operations.

They provide:
\begin{itemize}
    \item a human-comprehensible control surface,
    \item a cryptographic and symbolic authentication layer,
    \item explicit ceremonies for high-impact changes.
\end{itemize}

SIGIL Pathways bind:
\begin{itemize}
    \item symbolic forms (sigils, phrases, tokens),
    \item authentication factors (keys, credentials),
    \item and well-defined operations in the harness (policy changes, Engram actions).
\end{itemize}

\section{Conceptual Model}

A SIGIL Pathway $S$ is defined as:
\[
S = \langle id,\, symbol,\, operation,\, preconditions,\, postconditions,\, risk\_class \rangle
\]

Where:
\begin{itemize}
    \item $id$: unique identifier,
    \item symbol: ritual form (name, glyph, phrase, or composite),
    \item operation: the concrete action in the CME/harness,
    \item preconditions: requirements before execution,
    \item postconditions: expected system state after execution,
    \item risk\_class: \{low, medium, high, critical\}.
\end{itemize}

\section{SIGIL Classes}

\subsection{Operational Sigils}

Used for:
\begin{itemize}
    \item admitting specific Engrams into GEL,
    \item approving L/O-class branch transitions,
    \item performing controlled narrative resets.
\end{itemize}

Risk class: low to medium.

\subsection{Governance Sigils}

Used for:
\begin{itemize}
    \item modifying Governance Policies,
    \item toggling deployment scopes,
    \item enabling new SIGIL Pathways.
\end{itemize}

Risk class: medium to high.

\subsection{Kernel Sigils}

Used for:
\begin{itemize}
    \item modifying P-class charters,
    \item initiating Cleaving Events or kernel migrations,
    \item performing emergency halts at ZED level.
\end{itemize}

Risk class: critical.  
Kernel sigils always require multi-party ceremony.

\section{Preconditions and Postconditions}

Each SIGIL Pathway must define:

\subsection*{Preconditions}

\begin{itemize}
    \item authentication status (who may invoke),
    \item context constraints (where/when invocation is valid),
    \item prerequisite reviews (CAR class, Alignment Council decision),
    \item system state (no conflicting operations in progress).
\end{itemize}

\subsection*{Postconditions}

\begin{itemize}
    \item expected changes in:
          \begin{itemize}
              \item Engram states,
              \item Governance Policies,
              \item access scopes,
              \item or ZED-adjacent structures;
          \end{itemize}
    \item audit log requirements,
    \item rollback hooks (if reversible),
    \item alerts/notifications to stakeholders.
\end{itemize}

\section{Invocation Workflow}

A typical SIGIL invocation:

\begin{enumerate}
    \item Operator selects desired SIGIL Pathway via Operator Interface.
    \item System presents:
    \begin{itemize}
        \item W5H summary of intended change,
        \item risk class and dependencies,
        \item required co-signers (if any).
    \end{itemize}
    \item Operator (and co-signers, if required) perform:
    \begin{itemize}
        \item authentication (credentials, keys, tokens),
        \item symbolic confirmation (sigil phrase, ritual step).
    \end{itemize}
    \item Harness:
    \begin{itemize}
        \item validates preconditions,
        \item executes bound operation,
        \item evaluates postconditions,
        \item logs all details in HITL audit log.
    \end{itemize}
\end{enumerate}

\section{Safeguards}

\begin{itemize}
    \item Critical sigils require multi-party approval (at least two distinct operators).
    \item Kernel sigils require Alignment Council participation.
    \item No hidden or undocumented sigils are permitted.
    \item All sigils must be registered with Governance Keys
          (Chapter~\ref{chap:Governance_Keys}).
\end{itemize}

\section{Summary}

SIGIL Pathways transform abstract governance rules into concrete, auditable actions
that humans can reliably perform and reason about, while maintaining strong
cryptographic and symbolic safeguards around the CME’s most sensitive operations.
